% Please do not change %
\documentclass[11pt]{article} % font size
\usepackage[margin=1in]{geometry} % margins
\usepackage{amsmath,amsthm,amssymb} % for the  common "math symbols"
\usepackage{algorithm}
\usepackage{algpseudocode}
\usepackage[shortlabels]{enumitem} % for enumeration
\usepackage{booktabs} % if you need tables
\usepackage{listings}
\usepackage{color}

\lstset{
  frame=none,
  xleftmargin=2pt,
  stepnumber=1,
  numbers=left,
  numbersep=5pt,
  numberstyle=\ttfamily\tiny\color[gray]{0.3},
  belowcaptionskip=\bigskipamount,
  captionpos=b,
  escapeinside={*'}{'*},
  language=haskell,
  tabsize=2,
  emphstyle={\bf},
  commentstyle=\it,
  stringstyle=\mdseries\rmfamily,
  showspaces=false,
  keywordstyle=\bfseries\rmfamily,
  columns=flexible,
  basicstyle=\small\sffamily,
  showstringspaces=false,
  morecomment=[l]\%,
}
% Feel free to include additional packages (if needed), but make sure it does not override the margin/font requirements.
%
\setlength{\parindent}{0pt} % no indent for new paragraphs.
\setlength{\parskip}{5pt} % distance between paragraphs, which will be used when you have a blank line in your LaTeX code.


%%%%%%%%%%%%%%%%%%%%%%%%%%%%%%%%%%%%%%%%%%%%%%%%%%%%%%%%%%%%%%%%%%%%%%%%%%%%%%
\title{Probability and Computing: \\Exercise Sheet 5}
%
\author{Wira}
\date{13 Nov 2023}
%
\begin{document}
%
\maketitle

%You may use the usual \LaTeX\ environments to structure your answers:
%\begin{align}
%        \label{basegnn}
%        \mathbf{h}_u^{(t+1)} = \sigma \bigg( \mathbf{W}_\text{self}^{(t)} \mathbf{h}_u^{(t)}  + \mathbf{W}_\text{neigh}^{(t)} \sum_{v \in N(u)} \mathbf{h}_v^{(t)} + \mathbf{b}^{(t)} \bigg). 
%\end{align}


%\begin{table}[h]
%\caption{A simple table.}
%\centering
%\begin{tabular}[t]{lrrrr}
%\toprule
%Graphs & Can & Be  & Fun & Too \\ 
%\midrule 
%$G_1$ & $1$  & $2$ & $3$ & $4$ \\ 
%$G_2$ & $-1$  & $-2$ & $-3$ & $-4$ \\ 
%$G_3$ & $0.1$  & $0.2$ & $0.3$ & $0.4$ \\ 
%\bottomrule
%\end{tabular}
%\label{tab}
%\end{table}

\section*{Question 1}

% To answer each part of the question, please simply use a paragraph command 
% \paragraph{(a)}
Suppose we have $n$ variables and $t$ clauses. 
The variables within each clause are assumed to be distinct.

Let $m_i$ be the number of variables required to be true, 
and $l_i$ be the total number of variables in clause $C_i$.
Clauses with $m_i > k$ variables that need to be true 
will not have satisfying assignments with exactly $k$ true variables.
Thus, we can ignore them and only look at clauses that need $m_i \leq k$ variables to be true.

By Theorem 4.3.5, we can use the union-of-sets sampling algorithm to create an FPRAS solution.
The set $U$ is all $2^n$ truth assignments. 
The size of the input is $t + \log |U| = t + n$.

For $1 \leq i \leq t$, the set $H_i$ is the set of truth assignments 
with exactly $k$ true variables that satisfy clause $C_i$.
Then,
\[|H_i| = \binom{n-p_i}{k-m_i}\]
as this corresponds to setting $m_i$ variables true, 
and choosing which of the remaining $(n-p_i)$ variables not in clause $C_i$ to set to be true,
which can be computed in $poly(n)$ time.

To sample uniformly from $H_i$, we can pick $(k-m_i)$ variables without replacement u.a.r.
from the $(n-p_i)$ variables not in clause $C_i$.

Finally, given a truth assignment $x\in U$, we can check in poly time whether $x \in H_i$.

\clearpage
\section*{Question 2}

By Theorem 4.3.5, we can use the union-of-sets sampling algorithm to create an FPRAS solution.
The set $U$ is all $n!$ permutations.
The size of the input is $t+\log |U| = t+n\log n$.

For $1 \leq i \leq t$, the set $H_i$ is the set of permutations
that are increasing on set $S_i$.
Then,
\[|H_i| = \binom{n}{k_i}(n-k_i)!=\frac{n!}{k_i!}\]
since it is the same as choosing $k_i$ elements from the set $\{1,...,n\}$, 
where only 1 of the $k_i!$ permutations of $k_i$ elements will be increasing,
multipled by the rest of the $(n-k_i)!$ permutations.
This can be computed in $poly(n)$ time.

To sample uniformly from $H_i$, we can pick $(n-k_i)$ elements without replacement from $\{1,...,n\}$ u.a.r.
The remaining $k_i$ elements will be in sorted increasing order following the indices decided by $S_i$.

Finally, given a permutation $x\in U$, we can check in poly time whether $x\in H_i$.

\clearpage
\section*{Question 3}

\paragraph{(a)}
Provided are 2 algorithms to part (a)

\begin{algorithm}
\caption{$A'(\varepsilon', \delta')$}\label{alg:cap}
\begin{algorithmic}
\State $x' \gets 0$
\For{$i=1...m$}
\State $x' \gets x' + A_i (\varepsilon',1-(1-\delta')^{1/m})$
\EndFor
\State \Return $x'$
\end{algorithmic}
\end{algorithm}

\begin{algorithm}
\caption{$A'(\varepsilon', \delta')$}\label{alg:cap}
\begin{algorithmic}
\State $x' \gets 0$
\For{$i=1...m$}
\State $x' \gets x' + A_i (\varepsilon'/m,\delta'/m)$
\EndFor
\State \Return $x'$
\end{algorithmic}
\end{algorithm}

\paragraph{Algorithm 1}
 
\[A'(\varepsilon', \delta') = \sum_{i=1}^m A_i (\varepsilon',1-(1-\delta')^{1/m})\]

Note that 
\begin{equation}
\begin{aligned}
\left(\forall i, 1\leq i \leq m, (1-\varepsilon')x_i 
\leq A_i(\varepsilon', 1-(1-\delta')^{1/m}) \leq (1+\varepsilon') x_i\right)\\
\implies
\left((1-\varepsilon')\sum_{i=1}^m x_i \leq x' \leq (1+\varepsilon)\sum_{i=1}^m x_i\right)
\end{aligned}
\end{equation}


\begin{equation}
\begin{aligned}
    \Pr \left[x(1-\varepsilon') \leq x' \leq x(1+\varepsilon')\right]
    &= \Pr \left[(1-\varepsilon')\sum_{i=1}^m x_i \leq x' \leq (1+\varepsilon)\sum_{i=1}^m x_i\right] \\
    &\geq \prod_{i=1}^m \Pr \left[(1-\varepsilon')x_i 
    \leq A_i(\varepsilon', 1-(1-\delta')^{1/m}) \leq (1+\varepsilon') x_i\right] \\
    &= \prod_{i=1}^m \Pr \left[1-\varepsilon' 
    \leq \frac{A_i(\varepsilon', 1-(1-\delta')^{1/m})}{x_i} \leq 1+\varepsilon'\right]\\
    &\quad\quad\text{(by independence)}\\
    &= \left((1-\delta')^{1/m}\right)^m\\
    &= 1-\delta'
\end{aligned}
\end{equation}

\clearpage
\paragraph{Algorithm 2}

Note that 
\begin{equation}
\begin{aligned}
\left(|x'-x|\geq \varepsilon'x'\right) \implies
\left(\exists i : \left|A_i\left(\frac{\varepsilon'}{m},\frac{\delta'}{m}\right)
-x_i\right|\geq\frac{\varepsilon'x'}{m}\right)
\end{aligned}
\end{equation}

\begin{equation}
\begin{aligned}
    \Pr \left[x(1-\varepsilon') \leq x' \leq x(1+\varepsilon')\right]
    &= 1 - \Pr \left[|x'-x|\geq \varepsilon'x'\right] \\
    &\geq 1 - \Pr \left[\exists i : 
    \left|A_i\left(\frac{\varepsilon'}{m},\frac{\delta'}{m}\right) - x_i \right| \geq \frac{\varepsilon'x_i}{m}\right] \\
    &\geq 1 - \sum_{i=1}^m \Pr \left[ 
    \left|A_i\left(\frac{\varepsilon'}{m},\frac{\delta'}{m}\right) - x_i \right| \geq \frac{\varepsilon'x_i}{m}\right] 
    \quad\text{(by union bound)}\\
    &\geq 1 - m\left(\frac{\delta'}{m}\right)\\
    &= 1-\delta'\\
    &= 1-\delta'
\end{aligned}
\end{equation}

\paragraph{(b)}

\begin{algorithm}
\caption{$A_i(\varepsilon,\delta)$}\label{alg:cap}
\begin{algorithmic}
\State $X \gets 0$
\For{$j=1,...,k$}
\State $b_j \gets B_i$
\If{$b_j\in H_i$}
\State $X \gets X+1$
\EndIf
\EndFor
\State \Return $X/k$
\end{algorithmic}
\end{algorithm}

Let $X_1,X_2,...,X_k$ be the i.i.d. indicator variables that equal 1 if $b_j \in H_i$, 
where $b_j$ is obtained via $B_i$, and $\mu = \mathbb{E}[X_i] = |H_i|/|U_i| = x_i$.
Then from Theorem 4.3.1, if $k\geq \frac{3\log\frac{2}{\delta}}{\varepsilon^2 x_i}$
\[\Pr (|A_i(\varepsilon,\delta) - x_i| \leq \varepsilon x_i ) \geq 1-\delta\]

Thus,
\[\Pr \left[1-\varepsilon \leq \frac{A_i(\varepsilon,\delta)}{x_i} \leq 1+\varepsilon\right] \geq 1-\delta\]

\clearpage
\paragraph{(c)}

\begin{algorithm}
\caption{$A'(\varepsilon',\delta')$}\label{alg:cap}
\begin{algorithmic}
\State $X \gets 1$\
\For{$i=1,...,m$}
\State $b_i \gets A_i(\varepsilon'/(2m), \delta'/m)$
\State $X \gets X \times b_i$
\EndFor
\State \Return $X$
\end{algorithmic}
\end{algorithm}

This results in
\[A'(\varepsilon',\delta') = \prod_{i=1}^m A_i\left(\frac{\varepsilon'}{2m}, \frac{\delta'}{m}\right)\]

From part (b), we have
\[\Pr \left[\left|A_i\left(\frac{\varepsilon'}{2m}, \frac{\delta'}{m}\right) - x_i\right| 
\geq \frac{\varepsilon' x_i}{2m} \right]
\leq \frac{\delta'}{m}\]

Then the probability that all estimates have 
$\left|A_i\left(\frac{\varepsilon'}{2m}, \frac{\delta'}{m}\right) - x_i\right| 
\leq \frac{\varepsilon' x_i}{2m}$ 
is
\begin{equation}
\begin{aligned}
    \Pr \left[\forall i : \left|A_i\left(\frac{\varepsilon'}{2m}, \frac{\delta'}{m}\right) - x_i\right| 
    \leq \frac{\varepsilon' x_i}{2m} \right]
    &= 1 - \Pr \left[\exists i : \left|A_i\left(\frac{\varepsilon'}{2m}, \frac{\delta'}{m}\right) - x_i\right| 
    \geq \frac{\varepsilon' x_i}{2m} \right]\\
    &\geq 1 - \sum_{i=1}^m \Pr \left[\left|A_i\left(\frac{\varepsilon'}{2m}, \frac{\delta'}{m}\right) - x_i\right| 
    \geq \frac{\varepsilon' x_i}{2m} \right]
    \quad\text{(by union bound)}\\
    &\geq 1 - m\left(\frac{\delta'}{m}\right)\\
    &= 1 - \delta'
\end{aligned}
\end{equation}

From the proof of Claim 1 in Theorem 4.4.5, we then have, using 
\[\forall\varepsilon \in (0,1),1+\varepsilon\leq e^\varepsilon, 
1-\varepsilon/2\geq e^{-\varepsilon}, e^{\varepsilon/2}\leq 1 + \varepsilon,\]

\[1-\varepsilon' \leq e^{-\varepsilon'} \leq \left(1 - \frac{\varepsilon'}{2m}\right)^m 
\leq \prod_{i=1}^m \frac{A_i\left(\frac{\varepsilon'}{2m}, \frac{\delta'}{m}\right)}{x_i} 
= \frac{A'(\varepsilon',\delta')}{x}
\leq \left(1 - \frac{\varepsilon'}{2m}\right)^m 
\leq e^{\frac{\varepsilon'}{2}} \leq 1 + \varepsilon' .\]

Thus,
\[\Pr \left[1-\varepsilon' \leq \frac{A'(\varepsilon',\delta')}{x} \leq 1+\varepsilon'\right] \geq 1 - \delta'\]

From part (b) and Theorem 4.3.1, each $B_i$ is called at least 
\[k\geq \frac{3\log\frac{2}{\delta'/m}}{(\varepsilon'/2m)^2 x_i} 
= O\left(\frac{m^2 \log\frac{m}{\delta'}}{\varepsilon'^2}\right)\]
times.
\end{document}
